\documentclass[11pt]{article}

\usepackage[T1]{fontenc}
\usepackage[margin=0.85in]{geometry}
\usepackage{amsmath,amssymb,amsthm,mathtools,mathrsfs}
\IfFileExists{newtxtext.sty}{%
  \usepackage{newtxtext,newtxmath}%
}{%
  \usepackage{mathptmx}%
}
\usepackage{microtype}
\usepackage{booktabs}
\usepackage{enumitem}
\usepackage{xcolor}
\usepackage[
  colorlinks=true,
  linkcolor=blue!45!black,
  urlcolor=blue!45!black
]{hyperref}

\allowdisplaybreaks
\setlist[enumerate]{leftmargin=1.8em,itemsep=0.25em,topsep=0.4em}
\newtheorem{theorem}{Theorem}[section]
\newtheorem{proposition}[theorem]{Proposition}
\newtheorem{corollary}[theorem]{Corollary}
\theoremstyle{remark}
\newtheorem{remark}[theorem]{Remark}

\newcommand{\R}{\mathbb R}
\newcommand{\Risk}{\mathcal R}
\newcommand{\Transfer}{\mathcal T}
\newcommand{\silver}{\mathrm{sil}}

\title{\bfseries
Predetermined-to-Anytime Transfer for Gradient Descent\\
\large A Suffix-Local Addendum}
\author{Zuang Wang and Dandan Wang}
\date{August 20, 2026}

\begin{document}
\maketitle

\begin{abstract}
We record the precise transfer property connecting two lower-bound regimes
for last-iterate gradient descent on smooth convex functions.  By class
inclusion alone, a generic finite-horizon minimax ceiling transfers to an
anytime ceiling at the same exponent.  The stronger map
\[
 \Transfer(p)=\frac{2p}{1+p}
\]
requires more information: the same finite-horizon hard-instance certificate
must localize to every relevant ranked suffix at record prefixes.  We give a
self-contained suffix-local compiler theorem and its near/far proof.  The
currently audited joint-router certificate satisfies the compiler interfaces,
so one computer-assisted package proves simultaneously
\[
 \boxed{
 p_{\rm minimax}\le\frac{165737}{100000}=1.65737,
 \qquad
 q_{\rm any}\le\frac{331474}{265737}
 =1.2473761651557743\ldots .}
\]
These are supremal-exponent ceilings.  We do not claim endpoint attainment,
a positive uniform constant at either endpoint, or an endpoint
\(\mathrm{not}\text{-}o\) statement.
\end{abstract}

\section{Two schedule models}

Let \(\Risk_n(h_0,\ldots,h_{n-1})\) be the normalized worst-case
last-iterate objective gap of gradient descent over convex functions with
one-Lipschitz gradient.  All stepsizes below are nonnegative and fixed before
oracle information is observed.

In the predetermined-horizon model, the horizon is announced first and a
different finite schedule may be chosen for each horizon.  Define
\[
 p_{\rm minimax}:=\sup\left\{p\ge0:
 \begin{array}{l}
 \text{there are \(C<\infty\) and schedules \(h^{(T)}\) such that}\\[-1mm]
 \Risk_T(h^{(T)})\le C(T+1)^{-p}\text{ for every \(T\ge1\)}
 \end{array}\right\}.
\]

In the anytime model, one infinite schedule must work at every sufficiently
late deterministic prefix:
\[
 q_{\rm any}:=\sup\left\{q\ge0:
 \begin{array}{l}
 \text{there are one infinite schedule \(h\), \(C<\infty\), and \(n_0\)}\\[-1mm]
 \text{such that \(\Risk_n(h_0,\ldots,h_{n-1})\le C(n+1)^{-q}\)}\\[-1mm]
 \text{for every \(n\ge n_0\)}
 \end{array}\right\}.
\]
The anytime class is more constrained because its finite prefixes must be
compatible with one another.

\begin{proposition}[Black-box prefix transfer]
\label{prop:black-box}
If a universal finite-horizon theorem proves
\(\Risk_T(g)\ge c(T+1)^{-p}\) for every length-\(T\) schedule \(g\), then
every infinite schedule obeys the same lower bound at every prefix.  Hence
\(q_{\rm any}\le p\).
\end{proposition}

\begin{proof}
Apply the finite-horizon theorem to the first \(T\) steps of the infinite
schedule.  No compatibility beyond taking a prefix is needed.
\end{proof}

Proposition~\ref{prop:black-box} does \emph{not} yield the smaller exponent
\(2p/(1+p)\).  That stronger conclusion uses internal, suffix-local data from
the hard-instance proof.

\section{The suffix-local transfer compiler}

Fix \(0<a<1\) and put
\begin{equation}
 p:=1+a,\qquad b:=\frac2a,\qquad
 \Theta_*:=\frac1{a(a+2)}=\frac1{p^2-1},\qquad
 \Transfer(p):=\frac{2p}{1+p}=\frac{2(1+a)}{2+a}.
 \label{eq:parameters}
\end{equation}

For a finite schedule, write its positive excesses in decreasing order as
\(c_1\ge\cdots\ge c_r>0\), and define
\begin{equation}
 B:=1+\sum_{t<n}\min\{h_t,1\},\qquad
 D_j:=B+\sum_{i=j+1}^r c_i.
 \label{eq:ranked-state}
\end{equation}
Thus \(D_0=1+\sum_{t<n}h_t\), \(D_r=B\), \(B\le n+1\), and
\(r\le n\).

\begin{theorem}[Suffix-local transfer]
\label{thm:compiler}
Suppose one fixed chronological ranked-path certificate has the following
properties uniformly at every deterministic prefix.

\begin{enumerate}
\item The empty-chain and strict-record hard instances give
\begin{equation}
 \Risk_n\ge\frac1{2D_0-1},
 \qquad
 c_1\le D_0\sqrt{\Risk_n}
 \label{eq:record-arms}
\end{equation}
at every prefix whose last step is a strict record exceeding one.

\item For every \(2\le q\le r\), all top-\(q\) excesses and all
\(q-1\) chronological edges occur in a hard-instance bound
\begin{equation}
 \mathcal C_n(h)\ge\frac{c_{\rm J}}{D_q}
 \exp\left\{\mu(z+q-S_v)
 +\sum_{i=1}^{q-1}\Delta(v_i,v_{i+1})\right\},
 \label{eq:chronological-path}
\end{equation}
where \(c_{\rm J}>0\), \(0<\mu\le1\), \(z\ge0\),
\(S_v=\sum_i v_i\), and \((v_i)\) is a chronological permutation of
\[
 \widehat v_i=\frac{2D_q}{qc_i},\qquad 1\le i\le q.
\]
The normalized worst-case risk satisfies \(\Risk_n(h)\ge\mathcal C_n(h)\).

\item Fix finite thresholds \(0<t_0<\cdots<t_{M-1}<\infty\), and set
\[
 \alpha_\ell:=\frac1q\#\{i:v_i\le t_\ell\},\qquad
 k_\ell:=q\alpha_\ell,
\]
\[
 \mathscr R_\ell(\alpha):=1+
 2\sum_{u=\ell}^{M-2}\frac{\alpha_{u+1}-\alpha_u}{t_u}
 +\frac{2(1-\alpha_{M-1})}{t_{M-1}}.
\]
The exact denominator router, including zero occupancy, satisfies
\[
 \frac{D_{k_\ell}}{D_q}\le\mathscr R_\ell(\alpha).
\]
Every state row with positive dual weight is either a cumulative occupancy
monotonicity row or a globally valid affine tangent row
\[
 \mathscr R_{\ell(m)}(\alpha)\ge T_m(\alpha_{\ell(m)}),
 \qquad T_m(x)\le x^{-a}\quad(0<x\le1).
\]
There is no positive-weight mean row or other state row, and
\[
 H:=\max_m\sup_{0\le x\le1}T_m(x)<\infty.
\]

\item There are finite \(Q\) and \(B_0\) such that every supported state
with \(q\ge Q\) satisfies the signed finite-rank estimate
\[
 \frac1q\sum_{i=1}^{q-1}\Delta(v_i,v_{i+1})\ge B_0.
\]
This estimate includes all physical edges, endpoint corrections, and the
unbounded state group.  Its endpoint reserve
\[
 \mathfrak r:=B_0-\mu(2\Theta_*-1)
\]
is positive.  There is a finite localization factor \(\Lambda>1\) with
\begin{equation}
 \Lambda^a\ge H,
 \qquad
 \delta:=\mathfrak r-\frac{\mu b}{\Lambda}>0.
 \label{eq:localization-reserve}
\end{equation}
\end{enumerate}

Then no infinite predetermined schedule satisfies
\(\Risk_n=O(n^{-s})\) for any \(s>\Transfer(p)\).  Consequently,
\begin{equation}
 \boxed{q_{\rm any}\le\Transfer(p)=\frac{2p}{1+p}.}
 \label{eq:transfer}
\end{equation}
\end{theorem}

\begin{proof}
Assume that one infinite schedule satisfies
\(\Risk_n\le A n^{-s}\) for some \(s>\Transfer(p)>1\).  The empty-chain
arm rules out bounded steps, so there are arbitrarily late strict-record
prefixes.  At such a prefix write \(\epsilon=\Risk_n\), choose
\[
 \eta=\frac1{4\Lambda},
 \qquad
 k=\left\lfloor\frac{\eta}{\sqrt\epsilon}\right\rfloor,
\]
and let \(q\) minimize \(P_j:=D_jj^a\) over \(k\le j\le r\).
For late records, \(k\ge\eta/(2\sqrt\epsilon)\).  The record arms imply
\begin{equation}
 D_m\ge(1-\eta)D_0\ge\frac{1-\eta}{2\epsilon}
 \qquad(0\le m\le\min\{k,r\}).
 \label{eq:record-mass}
\end{equation}
In particular, eventually \(r\ge k\).
Indeed, if \(r<k\), then \eqref{eq:record-mass} with \(m=r\) gives
\(B=D_r\ge c/\epsilon\), contradicting \(B\le n+1\),
\(\epsilon\le An^{-s}\), and \(s>1\).

If \(q<\Lambda k\), then \(D_q>3D_0/4\), while suffix minimality and
\(B\le n+1\), \(r\le n\) give
\[
 c\epsilon^{-(p+1)/2}
 \le D_kk^a
 \le\frac43D_qq^a
 \le\frac43Br^a
 \le\frac43(n+1)^p.
\]
Therefore
\[
 \epsilon\ge c'(n+1)^{-2p/(1+p)},
\]
contradicting \(s>\Transfer(p)\).

Suppose instead that \(q\ge\Lambda k\), and put \(\tau=k/q\le1/\Lambda\).
Suffix minimality first gives the boundary floor
\(0<\widehat v_i\le b\).  It also validates every supported tangent row.
Indeed, if \(\alpha_\ell\ge\tau\), then \(k_\ell\ge k\) and
\[
 \mathscr R_\ell(\alpha)\ge\frac{D_{k_\ell}}{D_q}
 \ge\left(\frac q{k_\ell}\right)^a
 =\alpha_\ell^{-a}\ge T_m(\alpha_\ell).
\]
If \(\alpha_\ell<\tau\), then \(k_\ell<k\) and
\[
 \mathscr R_\ell(\alpha)\ge\frac{D_{k_\ell}}{D_q}
 \ge\frac{D_k}{D_q}\ge\left(\frac qk\right)^a
 \ge\Lambda^a\ge H\ge T_m(\alpha_\ell).
\]
The latter argument includes \(\alpha_\ell=0\), so no \(0^{-a}\) is
evaluated.  Suffix minimality also yields weak reciprocal supermajorization
against \(b((i-1/2)/q)^p\).  Convex midpoint quadrature then controls the
actual full-path mean without deleting a vertex or edge:
\begin{equation}
 S_v\le q\left(2\Theta_*+\frac b\Lambda\right).
 \label{eq:physical-mean}
\end{equation}
Since \(q\ge\Lambda k\ge c\epsilon^{-1/2}\ge c'n^{s/2}\), eventually
\(q\ge Q\).  Item 4 and equations \eqref{eq:chronological-path},
\eqref{eq:localization-reserve}, and \eqref{eq:physical-mean} imply
\begin{equation}
 \mathcal C_n(h)\ge c_{\rm J}\frac{e^{\delta q}}{D_q}.
 \label{eq:far-arm}
\end{equation}
Since \(D_qq^a\le Br^a\le(n+1)^p\), the assumed risk upper bound and
\eqref{eq:far-arm} force \(q=O(\log n)\).  On the other hand,
\(q\ge\Lambda k\ge c\epsilon^{-1/2}\ge c'n^{s/2}\), a contradiction.
The near and far branches prove \eqref{eq:transfer}.
\end{proof}

\section{Current exact instantiation}

\begin{corollary}[Audited paired ceilings]
\label{cor:instantiation}
The independently audited joint-router certificate at
\[
 p=\frac{165737}{100000},\qquad
 a=\frac{65737}{100000},\qquad
 \mu=\frac{8161}{10000}
\]
satisfies every hypothesis of Theorem~\ref{thm:compiler}.  It has
\[
 Q=10^9,\qquad
 B_0=\frac{118264177243}{10^{12}},\qquad
 H=\frac{63159358713716857163}{1250000000000000000}<60.
\]
Taking \(\Lambda=10^6\), the remaining exact suffix reserve is
\[
 \delta=
 \frac{133808620593933067}
 {17468753169000000000000}>0.
\]
\begin{equation}
 \boxed{
 p_{\rm minimax}\le\frac{165737}{100000}=1.65737,
 \qquad
 q_{\rm any}\le\frac{331474}{265737}
 =1.2473761651557743\ldots .}
 \label{eq:paired-result}
\end{equation}
\end{corollary}

The known silver constructions give the complementary achievable exponents.
Writing \(s_\silver=\log_2(1+\sqrt2)\), the current audited research
intervals are therefore
\begin{equation}
 \boxed{
 s_\silver\le p_{\rm minimax}\le\frac{165737}{100000},
 \qquad
 \frac{2s_\silver}{1+s_\silver}
 \le q_{\rm any}\le\frac{331474}{265737}.}
 \label{eq:intervals}
\end{equation}
Numerically, the constructive anytime endpoint is
\(2s_\silver/(1+s_\silver)=1.119545204061651\ldots\).

\begin{table}[ht]
\centering
\begin{tabular}{@{}lll@{}}
\toprule
Regime & achievable exponent & universal ceiling\\
\midrule
Predetermined horizon
 & \(\log_2(1+\sqrt2)=1.2715533\ldots\)
 & \(165737/100000=1.65737\)\\
Anytime
 & \(1.119545204061651\ldots\)
 & \(331474/265737=1.2473761651\ldots\)\\
\bottomrule
\end{tabular}
\caption{Current independently audited research intervals.  Smaller universal
ceilings are stronger.}
\end{table}

\clearpage
\section{Scope and evidence}

Theorem~\ref{thm:compiler} transfers a \emph{certificate}, not merely a
number.  In particular, the implication
\[
 p_{\rm minimax}\le p
 \quad\Longrightarrow\quad
 q_{\rm any}\le\frac{2p}{1+p}
\]
is not valid as a black-box theorem.  A finite-horizon induction may compress
past the suffix selected at a stopping-time record and therefore fail to
provide the local inequality used above.  Within the suffix-local certificate
class the map is monotone,
\[
 \Transfer'(p)=\frac2{(1+p)^2}>0,
\]
so every certified improvement in \(p\) automatically improves the anytime
ceiling as well.

The exact p=1.65737 verifier, certificate, controlled-repair replay, generic
suffix-local compiler, and independent audits all pass.  The principal
independent audit is
\begin{center}
\small
\texttt{audit\_joint\_tilted\_common\_multiplier\_p165737\_independent\_20260820\_v1.md}.
\end{center}
Its clean isolated replay and all transitive manifests pass.  The audit
independently checks the 199 positive dual rows, all 162,735 pair costs,
exact stationarity, the signed potential, finite-\(Q\) corrections, the
global tangent caps, and the positive suffix reserve displayed above.

This addendum is a research promotion of the paired ceilings, not a rewrite
of the hash-pinned \(\sqrt3\) or \(607/366\) release reports.  It makes no
claim that either remaining interval is closed.

\end{document}
